\begin{tikzpicture}[x=0.72cm,y=2.55cm]
% ---- (a) equilibrium landscape, A=0 ----
\draw[->] (-0.2,-0.65) -- (6.4,-0.65) node[below,font=\scriptsize] {reaction coordinate};
\draw[->] (-0.2,-0.65) -- (-0.2,1.3) node[above,font=\scriptsize] {free energy};
\draw[thick] plot[smooth] coordinates {(0,0.55) (0.5,0.1406) (1,0.1) (1.5,0.2969) (2,0.6) (2.5,0.8781) (3,1.0) (3.5,0.8781) (4,0.6) (4.5,0.2969) (5,0.1) (5.5,0.1406) (6,0.55)};
\node[below,font=\scriptsize] at (1,0.08) {filled site};
\node[below,font=\scriptsize] at (5,0.08) {empty site$+$ion};
\node[above,font=\scriptsize] at (3,1.0) {transition state};
\draw[densely dotted,gray] (1,0.1) -- (2.0,0.1);
\draw[densely dotted,gray] (3,1.0) -- (2.0,1.0);
\draw[{Latex[length=1.2mm]}-{Latex[length=1.2mm]}] (2.0,0.11) -- (2.0,0.99) node[pos=0.45,left,font=\scriptsize] {$\Delta G_a$};
\node[font=\scriptsize] at (3,-0.55) {(a) equilibrium: $\mathcal A=0$, Eyring $k\simeq k_0e^{-\Delta G_a/RT}$};
\begin{scope}[xshift=8.3cm]
% ---- (b) driven landscape: A=0.6, chi=1/2 — anchors from eq:bv (T4_calc.py) ----
\draw[->] (-0.2,-0.65) -- (6.4,-0.65) node[below,font=\scriptsize] {reaction coordinate};
\draw[densely dotted,thin,gray] plot[smooth] coordinates {(0,0.55) (0.5,0.1406) (1,0.1) (1.5,0.2969) (2,0.6) (2.5,0.8781) (3,1.0) (3.5,0.8781) (4,0.6) (4.5,0.2969) (5,0.1) (5.5,0.1406) (6,0.55)};
\draw[thick] plot[smooth] coordinates {(0,0.55) (0.5,0.1594) (1,0.1) (1.5,0.2481) (2,0.48) (2.5,0.6719) (3,0.7) (3.5,0.4844) (4,0.12) (4.5,-0.2544) (5,-0.5) (5.5,-0.4781) (6,-0.05)};
\node[below,font=\scriptsize] at (1,0.08) {filled site};
\node[below,font=\scriptsize] at (5,-0.58) {empty site$+$ion};
% forward barrier: DeltaG_a - chi*A, from reactant well to TS
\draw[densely dotted,gray] (1,0.1) -- (1.7,0.1);
\draw[densely dotted,gray] (3,0.7) -- (1.7,0.7);
\draw[{Latex[length=1.2mm]}-{Latex[length=1.2mm]}] (1.7,0.11) -- (1.7,0.69) node[pos=0.5,left,font=\scriptsize,align=right] {$\Delta G_a{-}\chi\mathcal A$};
% reverse barrier: DeltaG_a + (1-chi)*A, from TS to product well
\draw[densely dotted,gray] (3,0.7) -- (4.3,0.7);
\draw[densely dotted,gray] (5,-0.5) -- (4.3,-0.5);
\draw[{Latex[length=1.2mm]}-{Latex[length=1.2mm]}] (4.3,0.69) -- (4.3,-0.49) node[pos=0.55,right,font=\scriptsize,align=left] {$\Delta G_a{+}(1{-}\chi)\mathcal A$};
% driving force A: well-to-well drop
\draw[densely dotted,gray] (1,0.1) -- (5.85,0.1);
\draw[densely dotted,gray] (5,-0.5) -- (5.85,-0.5);
\draw[{Latex[length=1.2mm]}-{Latex[length=1.2mm]}] (5.85,0.09) -- (5.85,-0.49) node[pos=0.5,right,font=\scriptsize] {$\mathcal A$};
\node[font=\scriptsize] at (3,-0.85) {(b) driven $\mathcal A_j{=}sF(V_n{-}U_j){>}0$, $\chi{=}1/2$};
\end{scope}
\end{tikzpicture}
\caption{활성화 장벽 그림(좌표는 다섯 지점 — 반응물$\cdot$전이상태$\cdot$생성물 웰의 정확 높이(eq.~\eqref{eq:bv}
로 고정, 검산 = T4\_calc.py) — 사이의 매끄러운 보간이며, 보간 자체는 물리식이 아니라는 점을 명시한다). (a)
평형($\mathcal A=0$) — 두 골짜기 같은 높이, 속도는 Eyring 식 $k\simeq k_0e^{-\Delta G_a/RT}$. (b) 구동력
$\mathcal A_j$ 가 도착 골짜기를 내려($\Delta$골짜기 낙차 $=\mathcal A$, 우측 화살), 정방향 장벽이
$\Delta G_a-\chi\mathcal A$(좌측 화살, $\chi$ 몫)로 낮아지고 역방향 장벽은 $\Delta G_a+(1-\chi)\mathcal A$(우측
화살, $1-\chi$ 몫)로 높아진다 — 두 값이 같은 전이상태 높이에서 일관됨을 T4\_calc.py 가 재확인한다. 점선은
(a)(비교용). 상호작용 $\Omega$ 가 있는 전이에서는 같은 분할이 유효 장벽
$\Delta H_{a,j}^\eff=\Delta H_{a,j}-\chi_d\Omega_j$(N7, 표~\ref{tab:staging} 열)로 일반화된다. 이 그림이
식~\eqref{eq:xieq} logistic 과 \S\ref{sec:lag} 의 $L_q$ 두 결과의 공통 출발점이다.}
\label{fig:barrier}
